\subsection{Ordinary algorithms over effective henselian discrete valuation rings}
\label{C:henselian-effective-subsection}

This subsection uses finite names for ring elements, rather than
infinite precision streams.  The hypotheses on those names are made
explicit, and are then proved for a noncomplete henselization.
The integral coefficient modules $C_d$ are those of
Lemma~\ref{C:integral-orthogonal-coalgebras}.

\begin{lemma}[Proved: henselian idempotent and hermitian lifting]
\label{C:henselian-finite-algebra-lifting}
Let $(R,\mathfrak m)$ be a henselian Noetherian local ring with
$2\in R^\times$, and let $E$ be an associative unital algebra
finite as an $R$-module, with central $R$-action.  Put
$J=\operatorname{rad}E$.
Then $E$ is semiperfect, and idempotents lift through every
two-sided ideal contained in $J$.

If $E$ has an $R$-linear involution $*$ and $a=a^*$ satisfies
$a-1\in J$, there is an element $b=b^*$ with
$b^2=a$ and $b-1\in J$.  Consequently nonsingular hermitian
forms on a fixed finite projective $E$-module, with either sign,
are classified up to isometry by their reductions modulo $J$.
In particular the radical reduction and transfer description
for a category of finite free rational representations over $R$
has the same residue division-algebra form factors as in the
complete case; completeness is unnecessary here.
\end{lemma}
\begin{proof}
First $\mathfrak mE\subseteq J$.  If $h\in\mathfrak mE$,
left multiplication by $1-uh$, for any $u\in E$, is the
identity modulo $\mathfrak m$ on the finite $R$-module $E$.
It is surjective by Nakayama and is injective because a
surjective endomorphism of a Noetherian module is injective.
Its inverse supplies a two-sided inverse to $1-uh$.
Thus the Jacobson criterion applies.  The algebra
$E/\mathfrak mE$ is finite-dimensional over $R/\mathfrak m$,
so $E/J$ is semisimple Artinian.

We need a unit-reflection observation.  If $A\subseteq E$
is a commutative $R$-subalgebra finite over $R$, then $E$ is
a finite faithful left $A$-module: its finite $R$-generating
set also generates it over $A$, and $A\hookrightarrow E$
is an inclusion of $A$-modules.  An element $c\in A$ which
is a unit in $E$ is a unit in $A$.  Otherwise a maximal
ideal $\mathfrak n\subset A$ contains $c$.  Localization
gives the nonzero finite module $E_{\mathfrak n}$, whereas
invertibility of $c$ gives
$cE_{\mathfrak n}=E_{\mathfrak n}$, contradicting Nakayama.
It follows that
\[
                 A\cap J\subseteq\operatorname{rad}A.       \tag{C.3.HD1}
\]
Indeed $1-xh$ is a unit in $E$, hence in $A$, for every
$h\in A\cap J$ and $x\in A$.

Let $K\subseteq J$ be a two-sided ideal and suppose
$v^2-v\in K$.  The commutative subalgebra $A=R[v]$ is
finite over $R$, by the determinant trick for an element
of a finite algebra.  A finite commutative algebra over a
henselian local ring is a finite product of henselian local
rings \cite[Tag 04GH, Lemma 10.153.4(1)]{C-StacksHenselianEffective}.
By (C.3.HD1), $K\cap A\subseteq\operatorname{rad}A$.
Each local factor of $A/(K\cap A)$ has only the idempotents
$0,1$.  Choose those same zeros and ones in the corresponding
local factors of $A$.  The resulting idempotent is congruent
to $v$ modulo $K\cap A$, hence modulo $K$ in $E$.
This proves the lifting assertion, not just lifting modulo
$\operatorname{rad}A$.  Lifting a complete set of primitive
orthogonal idempotents successively in complementary corners
shows that $E$ is semiperfect.

For the square-root assertion put $A=R[a]$.  This is again
finite and commutative, and the involution fixes it pointwise.
By (C.3.HD1), $a-1\in\operatorname{rad}A$.
In each henselian local factor of $A$, the polynomial
$T^2-a$ has the simple residual root $1$, since $2$ is a
unit.  Hensel's lemma supplies $b\in A$ with
$b^2=a$ and $b\equiv1\pmod{\operatorname{rad}A}$
\cite[Tag 04GG, Lemma 10.153.3(2)]{C-StacksHenselianEffective}.
Then $b+1$ is a unit in $A$, so
\[
                 b-1=(a-1)(b+1)^{-1}\in J.
\]
Also $b=b^*$, and $b$ is a unit.

Here is the form-lifting argument explicitly.  It applies
also in any finite matrix or corner algebra over $E$.
A nonsingular residue form on a given projective module
lifts by lifting its coefficients and symmetrizing with
$1/2$; a map invertible modulo the radical is invertible.
For injectivity, choose one form $q$, use it to define the
adjoint involution, and write another form as $q a$ with
$a=a^*$ a unit.  Lift a residue isometry to a unit $x$.
Then $d=x^*ax$ is self-adjoint and congruent to $1$
modulo the radical.  The square root just constructed
gives $b=b^*$, $b^2=d$; hence
$(xb^{-1})^*a(xb^{-1})=1$.
Thus the residue isometry lifts.

For representations the endomorphism rings are finite
$R$-modules, since they are submodules of finite matrix
modules and $R$ is Noetherian.  Their lifted idempotents
split into finite projective, hence free, $R$-modules.
Semiperfectness therefore gives Krull--Schmidt decompositions.
Apply the preceding form argument to the endomorphism ring
of each finite direct sum and its adjoint involution.
The resulting decomposition and transfer are precisely
the semisimple residue-block argument of
\cite[Theorem 2.2, Corollary 2.3, Propositions 2.4--2.5,
and Section 3, pp.\ 269--273]{C-QSS1979}.
The square-root argument above replaces the completeness
step in radical isometry lifting.  No assertion that $E$
is radical-complete has been made or used.
\end{proof}

Call a countable discrete valuation ring $R$ \emph{effective}
here if its elements have a computable enumeration without
repetition, arithmetic and equality are computable, the
valuation of a nonzero element is computable, and $a/b\in R$
is computable whenever $b\ne0$ and $v(a)\geq v(b)$.
Require also a supplied finite residue field table, a
computable reduction map, and computable lifts of its
elements.  These are exact finite-input requirements;
they do not include a decision oracle for infinite trees.

\begin{theorem}[Proved: complete ordinary classification over an effective henselian DVR]
\label{C:effective-henselian-DVR-WG}
Let $R$ be an effective henselian discrete valuation ring
with finite residue field of odd characteristic.
For every $m\geq0$, the indecomposable finite projective
rational representations of split $O_m$ have a computable
canonical index set $\mathscr I_R\subseteq\mathbf N$,
with no remaining isomorphism quotient.  Isomorphism,
indecomposability, duality, nonsingular invariant form
type and equivariant isometry are decidable from finite
coefficient data over $R$.

The self-dual indices have computable finite residue
endomorphism fields $D_i$, adjoint involutions $\sigma_i$
and signs $\varepsilon_i$.  With symmetric, alternating,
unitary and dual-pair index sets defined by
\[
 \begin{array}{ll}
 \mathscr I_{\rm s}:&\varepsilon_i=1,\quad\sigma_i=1,\\
 \mathscr I_{\rm a}:&\varepsilon_i=-1,\quad\sigma_i=1,\\
 \mathscr I_{\rm u}:&\sigma_i\ne1,\\
 \mathscr I_{\rm h}:&\text{the smaller index in each non-self-dual pair},
 \end{array}
\]
there are explicit actual form generators giving
\[
 WG_{\rm sp}(BO_{m,R})\simeq
 \bigoplus_{i\in\mathscr I_{\rm s}}
       (\mathbf Z e_i\oplus\mathbf Z/2\,t_i)
 \ \oplus\!
 \bigoplus_{i\in\mathscr I_{\rm a}\cup\mathscr I_{\rm u}
                           \cup\mathscr I_{\rm h}}
       \mathbf Z e_i.                                    \tag{C.3.HD2}
\]
All coordinates, products and specified exterior-power
coefficients are obtained by ordinary terminating algorithms.
The canonical set is countable; it is not asserted to be
finite at fixed rank and coefficient degree.
\end{theorem}
\begin{proof}
Every finite projective $R$-module is free.  An input given
as the image of an idempotent matrix can be put in a free
basis by selecting independent columns modulo the maximal
ideal; the unit minor and Nakayama prove that the selected
columns form a basis of the image.  Thus it suffices to
work with free modules and finite coefficient matrices
in the explicitly based modules $C_d(R)$.

\emph{Linear algebra and actual isomorphism tests.}
Kernels of finite matrices over $R$ are computable.
Choose a nonzero entry of least valuation as pivot.  It
divides every entry, so the permitted exact divisions and
invertible row and column operations clear its row and
column.  Repeat on the smaller matrix.  This diagonal
reduction terminates and supplies a basis of the kernel.
It is the usual Smith procedure over a discrete valuation
ring, with all division steps justified by the pivot choice.

Commutation of a constant matrix with two coactions is a
finite $R$-linear coefficient system.  Hence the full
$\operatorname{Hom}$ module, including a basis, is computable.
Two representations of the same free rank are isomorphic
if and only if the residue of some element of this Hom
module has nonzero determinant.  Its residues form the
finite vector space spanned by the reductions of the
computed generators.  Enumerate all that space, test
determinants in the finite residue field, and lift a
successful linear combination.  The lift has unit
determinant and is an actual equivariant isomorphism.
If the finite search fails no such isomorphism exists.

Compute $E=\operatorname{End}(M)$ in the same way, and let
$B$ be its image in
$\operatorname{End}_{\kappa}(M/\mathfrak mM)$.
This finite algebra is explicitly computable.  The kernel
of $E\to B$ lies in $J=\operatorname{rad}E$: if the matrix
of $h$ is zero modulo $\mathfrak m$, then $1-uh$ has unit
determinant for every $u\in E$, and its inverse matrix is
again equivariant.  Consequently
\[
                    E/J\simeq B/\operatorname{rad}B.       \tag{C.3.HD3}
\]
The radical of the finite ring $B$ is computed by testing,
for every $b\in B$, whether $1-xb$ is a unit for every
$x\in B$; inverses are found by finite enumeration.

By Lemma~\ref{C:henselian-finite-algebra-lifting}, $E$ is
semiperfect and its idempotents lift from $B$.
The module $M$ is indecomposable if and only if $B$ has
no nontrivial idempotent, a finite decidable condition.
If one exists, enumerate elements of the computed
finite-generator $R$-module $E$ until an idempotent with
that prescribed residue appears.  Equality is decidable,
and the lifting lemma proves that this search terminates.
Choose free bases for its image and complement and repeat
on the two resulting coactions.  Positive ranks decrease,
so this computes a Krull--Schmidt decomposition after at
most $\operatorname{rank}(M)-1$ splits.

\emph{Literal canonical indices.}
Fix computable natural-number codes for all finite tuples
of ring-element names, together with the rank and a
coefficient degree.  Coassociativity and the counit are
finite polynomial equations in the finite free coefficient
modules, hence are decidable.  Require minimal degree by
testing that at least one coefficient outside $C_{d-1}$
is nonzero, when $d>0$.  This condition is invariant under
a change of basis because $C_{d-1}$ is an $R$-submodule.
Also require indecomposability by the preceding finite test.
For a valid code $c$, keep it precisely when no smaller
valid code describes an isomorphic representation.
There are only finitely many smaller natural numbers,
and each required isomorphism test terminates.
This is a decidable subset $\mathscr I_R\subseteq\mathbf N$.
Every indecomposable has a finite code and therefore a
least such code.  Given any input code, comparison with
its finitely many predecessors computes its canonical
index.  Thus the construction removes, rather than
renames, the representation isomorphism problem.

\emph{Duality and invariant forms.}
The antipode and transpose compute the dual coaction,
and the same finite-predecessor procedure computes its
index.  For an indecomposable, (C.3.HD3) is a finite
division ring, hence a finite field
\cite[Theorem III.4.4, p.\ 29]{C-ArtinFiniteDivision}.
Symmetric and alternating invariant matrices are kernels
of finite linear systems over $R$.  The finite residue
determinant test used above decides whether such a space
contains a nonsingular matrix, and produces one if it does.
If the module is self-dual, at least one sign occurs:
for an isomorphism $g:M\to M^*$, the two endomorphisms
$g^{-1}(g+g^{\mathsf t})$ and
$g^{-1}(g-g^{\mathsf t})$ sum to $2$, and cannot both lie
in the radical of the local ring $E$.
Choose a symmetric form if possible, otherwise an
alternating form, and then the least valid matrix code
of that sign.  Denote it by $q_i$ and its sign by
$\varepsilon_i$.  Its adjoint involution on $E_i$ and
the induced $\sigma_i$ on $D_i$ are computable by matrix
operations and the finite quotient (C.3.HD3).

The lifting lemma and the finite-field hermitian
classification \cite[Lemma 3.1(a),(b)]{C-SergeichukFiniteForms}
give (C.3.HD2), with the same explicit generators as
in Theorem~\ref{C:complete-finite-WG}.
For clarity, in symmetric type put $e_i=[q_i]$ and
$t_i=[q_i a_i]-[q_i]$, where $a_i$ is an invertible
self-adjoint lift of a chosen nonsquare in $D_i$.
It is found by lifting a residue endomorphism through
the computed generators of $E_i$ and averaging with
its adjoint.  The class $t_i$ has exact order two by
the finite-field determinant invariant.
For alternating type and for a non-self-dual pair put
$e_i=H(M_i)$.  In unitary type choose a nonzero
$c\in D_i$ satisfying $\sigma_i(c)=\varepsilon_i c$;
take $c=1$ for positive sign and
$c=u-\sigma_i(u)\ne0$ for negative sign.
Lift and average to obtain $a_i^*=\varepsilon_i a_i$.
Then $e_i=[q_i a_i]$ is represented by a symmetric
nonsingular form.  These constructions take finite time.

\emph{Isometry and the coordinates of an actual form.}
First test whether its underlying module is isomorphic
to that of a second form.  After a successful isomorphism,
write the second form as $q a$ relative to the first,
and compute $E/J$ with its adjoint involution.
The forms are isometric if and only if some unit
$\bar x\in E/J$ satisfies
$\bar x^*\bar a\bar x=1$.
The finite quotient permits exhaustive testing.
The equivalence is exactly the radical isometry lifting
proved in the lemma.  If an actual isometry matrix is
requested, enumerate elements of $E$ and test the exact
equation $x^*ax=1$ until one occurs; the positive finite
test proves termination.  Hence no nonlinear equation
oracle over $R$ is being assumed.

The underlying Krull--Schmidt multiplicities of an
actual symmetric form leave only finitely many candidate
form coordinates: one determinant bit for each occurring
symmetric type and no additional bit in the other types.
A symmetric multiplicity $r$ with bit $b\in\{0,1\}$
has coordinate $r e_i+b t_i$ and representative consisting
of $r-b$ copies of $q_i$ and $b$ copies of $q_i a_i$.
Alternating multiplicity is even and contributes half
its value; unitary multiplicity contributes itself;
a dual pair contributes its common multiplicity.
Construct these finitely many candidate forms and use
the isometry test to select the unique coordinate list.
This also proves cancellation and decidable equality
in the displayed split group.

Finally tensor coactions multiply their coefficient
matrices, and exterior coactions and forms use minors.
The coefficient bounds are $d+e$ and $rd$, respectively.
Apply the terminating decomposition and coordinate
algorithms just proved.  Extend by bilinearity and by
$\lambda_t(-x)=\lambda_t(x)^{-1}$ to every specified
finite exterior coefficient on virtual inputs.
All data in this construction have finite names and
every search has the stated termination proof.
\end{proof}

\begin{theorem}[Proved: an explicit noncomplete henselian DVR family]
\label{C:function-DVR-henselization-effective}
Let $q$ be any odd prime power and put
\[
              R=\bigl(\mathbf F_q[t]_{(t)}\bigr)^h.
\]
There is a uniform exact effective presentation of $R$
satisfying every hypothesis of
Theorem~\ref{C:effective-henselian-DVR-WG}.
Consequently that theorem gives ordinary terminating
canonical classification, equality, products and
exterior operations for $WG_{\rm sp}(BO_{m,R})$ for
every $m$.  This $R$ is a positive-dimensional
Noetherian henselian local domain and is not complete.
No factorization or infinite-branch oracle is required.
\end{theorem}
\begin{proof}
Write $A_0=\mathbf F_q[t]_{(t)}$ and $K_0=\mathbf F_q(t)$.
The henselization is a discrete valuation ring, with
uniformizer $t$ and residue field $\mathbf F_q$
\cite[Tag 0AP3, Lemma 15.46.11]{C-StacksHenselianEffective}.
It has the same completion as $A_0$, namely
$\mathbf F_q[[t]]$
\cite[Tag 06LJ, Lemma 15.46.3]{C-StacksHenselianEffective},
and the separated discrete valuation embeds it in
that completion.

\emph{Finite names and computable expansions.}
Enumerate the data
\[
 \begin{gathered}
 (f,a,g,h),\qquad f\in A_0[T]\text{ monic},\quad
 a\in\mathbf F_q,\\
 \bar f(a)=0,\quad\bar f'(a)\ne0,\qquad
 g,h\in A_0[T],\quad\bar h(a)\ne0.
 \end{gathered}                                        \tag{C.3.HD4}
\]
They name $g(\alpha)/h(\alpha)$, where $\alpha$ is the
unique root of $f$ congruent to $a$ modulo $t$.
All validity conditions are decidable.  The root's
power-series coefficients are computable successively:
after a solution modulo $t^N$, the correction coefficient
at $t^N$ is obtained from a linear equation with the
nonzero coefficient $\bar f'(a)$.
Polynomial evaluation and inversion of $h(\alpha)$,
whose constant coefficient is nonzero, compute the
named element to every finite precision.

These data name every element of $R$.
The henselization is the filtered union of pointed
\'{e}tale neighborhoods with the same residue field
\cite[Tag 0BSK, Lemma 10.155.1 and its construction]
 {C-StacksHenselianEffective}.
Near the specified point, each neighborhood is standard
\'{e}tale, hence is a localization of $A_0[T]/(f)$
for a monic $f$, with $f'$ invertible at that point
\cite[Tag 00UE, Proposition 10.144.4]{C-StacksHenselianEffective}.
Its residue sends $T$ to some $a\in\mathbf F_q$.
Every element in this localization is therefore of
the form (C.3.HD4).  This proves coverage, including
elements initially described using several neighborhoods.

\emph{Factorization in the finite function fields used below.}
Polynomials over $K_0$ can be factored by a finite
elementary search.  Clear denominators and remove the
content to obtain $F\in\mathbf F_q[t,T]$.
If $D$ is its total degree, all polynomial factors have
total degree at most $D$.  There are only finitely many
normalized polynomials with that bound over $\mathbf F_q$.
Test divisibility and recursively split proper factors;
Gauss's lemma gives factorization over $K_0$.

Factorization of separable polynomials is also effective
over every finite extension supplied as a tower of
separable simple extensions of $K_0$.
Indeed the separable norm-and-shift induction in the
proof of Theorem~\ref{C:function-field-factorization}
uses only effective base factorization and an infinite
enumerable coefficient field.  Both hold here.
Explicitly, for $L=E[\gamma]$ separable and a monic
separable $u\in L[T]$, enumerate $c\in E$ until
$\operatorname{Norm}_{L/E}(u(T-c\gamma))$ is separable.
Only finitely many $c$ create collisions between the
finitely many conjugate roots, so this search terminates.
Factor that norm over $E$; the corresponding gcds after
extension to $L$ recover the shifted irreducible factors
of $u$.  Undo the shift.  The proof in the cited theorem
shows that a squarefree norm separates the factors
and gives multiplicity one; it applies in every
characteristic.  Induct on the tower.  All extensions
required below are separable, so no inseparable
factorization step is hidden in this use.

\emph{Exact fields behind the expansions.}
For a root datum $(f,a)$ in (C.3.HD4), factor $f$ over
$K_0$.  Its monic irreducible factors have coefficients
in $A_0$, since those coefficients are integral over
the integrally closed domain $A_0$.
Exactly one irreducible factor has $a$ as a root of its
reduction: the product has the simple residual root
$a$.  Select that factor, denoted $f_a$.
It is separable, since its derivative is nonzero at
the specified residual root.  Thus
$K_0[T]/(f_a)$ is an explicit finite separable field
with a specified embedding in $\mathbf F_q((t))$.
Its generator maps to the computable series $\alpha$.

To combine two finite sets of such roots, add the
roots one at a time.  Suppose that the current
explicit field $L$ already has its specified embedding
in $\mathbf F_q((t))$, and the next root is $\beta$,
with known separable minimal polynomial $f_b$ over
$K_0$.  Factor $f_b$ over $L$ by the preceding procedure,
as a product of distinct monic irreducible factors
$v_1,\ldots,v_s$.
Exactly one evaluation $v_j(\beta)$ is zero in
$\mathbf F_q((t))$.  There is at least one because
$f_b(\beta)=0$; there cannot be two because the
distinct factors are coprime and the embedding of
$L$ is injective.

Compute the Laurent expansions of these evaluations
in parallel.  Every nonzero one eventually exposes
a nonzero coefficient.  Stop when all but one have
done so; the remaining factor is the correct factor.
This is a terminating selection algorithm, without
assuming a zero test in the as-yet unknown compositum.
For $s=1$ it stops immediately.
Form $L[T]/(v_j)$ using the selected factor and send
its generator to $\beta$.  This is the required field
compositum inside $\mathbf F_q((t))$.

Every element of a current field is stored as a finite
expression in its original simple-root generators.
Its Laurent expansion is computable: exact arithmetic
already decides which denominators are nonzero in
that field, and expansion until a first nonzero
coefficient supplies their valuations and permits
Laurent inversion.  Thus the parallel evaluations
in the preceding paragraph are effective even when
their coefficients have poles.  Their finite
expressions give computable lower valuation bounds.
When two previously computed field elements are to
be compared, merge their finite lists of original
root generators by this procedure, rather than
assuming a new minimal-polynomial routine for an
arbitrary element.  Arithmetic and equality then
take place exactly in the explicit common field.

This constructs an exact equality test for all names
in (C.3.HD4), as well as exact field arithmetic for
finite expressions in them.  A nonzero expression's
valuation is computed by its first nonzero Laurent
coefficient.  The test for being zero is performed
first, so the valuation search is never applied to
zero.  Since $R$ is a discrete valuation ring, it is
the valuation ring of its fraction field.  Therefore
the quotient $x/y$ of two named elements with
$y\ne0$ and $v(x)\geq v(y)$ belongs to $R$.
It is already available as an exact field expression.
To give it a name of type (C.3.HD4), enumerate those
names and use the equality test until it is found.
Coverage proves termination.  The same procedure
returns such a name for a sum, product or additive
inverse.  Reduction modulo $t$ is its constant
series coefficient, and constant polynomials give
the residue-field lifts.

Finally remove duplicates from the effective
enumeration (C.3.HD4) using its exact equality test,
retaining the first occurrence of each element.
This gives a computable enumeration without
repetition.  It has infinitely many entries,
already because $A_0\subset R$; hence its $n$th
entry is found in finite time for every $n$.
An arithmetic result's index is found by the same
terminating equality search.  All required access
to $R$ has now been supplied constructively.

The ring is countable by its finite names, whereas
$\mathbf F_q[[t]]$ is uncountable.  It is therefore
not complete.  The discrete valuation assertion
already proves its Noetherian, local, domain and
positive-dimensional properties.  Applying
Theorem~\ref{C:effective-henselian-DVR-WG} completes
all the stated classification and operation claims.
\end{proof}

\begin{theorem}[Proved: every affine finite-field local henselization]
\label{C:finite-field-henselization-WG}
Let $k$ be a supplied finite field of odd characteristic,
let $S=k[z_1,\ldots,z_s]/I$ be a supplied finite
presentation, and let $S\twoheadrightarrow\kappa$ be
a supplied map onto a finite field.  Write
$\mathfrak m$ for its kernel and put $R=(S_{\mathfrak m})^h$.
Then $R$ has an explicit finite-name presentation with
computable arithmetic, equality, residue and unit
inversion, and there is an ordinary terminating
algorithm computing a finite generating set for the
kernel of every finite matrix over $R$.

For every $m\geq0$, all conclusions of
Theorem~\ref{C:effective-henselian-DVR-WG}, including
the literal computable indecomposable indices,
(C.3.HD2), isometry and equality decisions, products
and exterior operations, hold for this $R$ by ordinary
terminating algorithms.  No regularity, reducedness,
DVR, or positive-dimensional hypothesis is imposed
on $S_{\mathfrak m}$.  If $R$ is finite, its canonical
finite-name set is a decidable subset of $\mathbf N$;
no bijective enumeration by all of $\mathbf N$ is asserted.
\end{theorem}
\begin{proof}
We specify both the finite algebra computations and
how they compute the actual henselization rather
than a formal completion.

\emph{Finite presentations and their kernels.}
Over a finite field, polynomial ideal membership,
annihilators and generators of kernels of finite
polynomial matrices are computable by terminating
Gr\"obner algorithms.  The required elimination and
syzygy constructions are described in the primary
account \cite[Sections 4.3--4.4 and 5.1--5.3,
pp.\ 240--249]{C-BuchbergerSystems2001}.
Here is the reduction for an arbitrary presented
algebra $T=k[w_1,\ldots,w_a]/(F_1,\ldots,F_b)$.
For a matrix $A:T^u\to T^v$, choose a polynomial
lift $\widetilde A$.  Compute the polynomial
syzygies of the matrix whose columns are those
of $\widetilde A$ together with the $bv$ columns
$-F_j e_l$ for $1\leq j\leq b$, $1\leq l\leq v$.
Project each syzygy to its first $u$ coordinates
and reduce modulo $(F_1,\ldots,F_b)$.
These vectors generate $\ker A$: an element of
that kernel is exactly a polynomial vector whose
image under $\widetilde A$ is a combination of
the $F_j e_l$.
In particular the same construction for multiplication
by $a\in T$ computes $\operatorname{Ann}_T(a)$.
Localizations cause no difficulty: inverting $g$
adds a variable $y$ and the relation $yg-1$.
All rings used below have such finite presentations.

\emph{Finite pointed neighborhoods.}
A name consists of a finitely presented \'etale
$S$-algebra $T$, a specified extension
$\rho:T\twoheadrightarrow\kappa$ of the given
residue map, and an element of $T$ in polynomial
coordinates.  For effective enumeration it is
unnecessary to decide \'etaleness of arbitrary
presentations.  Enumerate the standard \'etale
ones as follows.  Choose $b\in S$ with
$\bar b\ne0$ in $\kappa$, a monic
$f\in S_b[U]$, a root $a\in\kappa$ of
$\bar f$ with $\bar f'(a)\ne0$, and a polynomial
$h\in S_b[U]$ with $\bar h(a)\ne0$.
Use
\[
 T=S_b[U,V]/\bigl(f(U),\,V f'(U)h(U)-1\bigr),
 \qquad
 \rho(U)=a,\qquad
 \rho(V)=\bigl(\bar f'(a)\bar h(a)\bigr)^{-1}.       \tag{C.3.HD5}
\]
The relation inverting $b$ is included in its
polynomial presentation over $k$.
The displayed residue tests and all its finite
polynomial data are decidable.  This construction
is \'etale and defines the indicated residue map.
Allow arbitrary elements of the resulting $T$.
For internal arithmetic also allow finite tensor
products of these algebras over $S$, with their
product residue map.  These are again pointed
finitely presented \'etale neighborhoods, with
explicit polynomial presentations and maps.

Every element of $R$ has such a name.
Indeed the pointed \'etale construction of
henselization and the standard \'etale neighborhood
lemma used in the preceding theorem apply to
$S_{\mathfrak m}$ as well.  The finitely many
coefficients in a standard presentation over
$S_{\mathfrak m}$ all lie in one $S_b$ with
$b\notin\mathfrak m$.  Shrinking by the polynomial
$h$ and by $f'$ gives exactly (C.3.HD5).
This also permits all the elements and denominators
of any finite collection of names to be included
in one such tensor neighborhood.

For any such $T$ put $\mathfrak q=\ker\rho$.
There is a canonical identification
\[
                  (T_{\mathfrak q})^h\simeq R.           \tag{C.3.HD6}
\]
For completeness, the local map
$S_{\mathfrak m}\to T_{\mathfrak q}$ is a localization
of an \'etale map and induces the identity on
$\kappa$.  The pointed map $T\to R$ exists and
is unique by the henselian lifting property.
The universal properties of the two henselizations
then give maps in both directions in (C.3.HD6).
Their compositions are identities: on $S_{\mathfrak m}$
this is uniqueness of its map to a henselian local
ring, and on $T$ it is uniqueness of the pointed
\'etale lift.  The same argument applies after
localizing at $\mathfrak q$.
Equivalently (C.3.HD6) is the cofinality of pointed
\'etale neighborhoods in the henselization construction
\cite[Section 10.155; also the proof of Lemma 15.46.3(c)]
 {C-StacksHenselianEffective}.
It follows that $T\to R$ is flat and that
$T_{\mathfrak q}\to R$ is faithfully flat
\cite[Lemma 15.46.1(1), Tag 07QL]{C-StacksHenselianEffective}.

\emph{Exact equality and arithmetic.}
For a named element $z\in T$, faithful flatness gives
\[
 \begin{split}
 z=0\text{ in }R
 &\ \Longleftrightarrow\ z=0\text{ in }T_{\mathfrak q}\\
 &\ \Longleftrightarrow\
       \operatorname{Ann}_T(z)\not\subseteq\mathfrak q.
 \end{split}                                             \tag{C.3.HD7}
\]
Compute generators of that annihilator by the finite
procedure above.  The last condition is equivalent
to at least one generator having nonzero image under
$\rho$, and is a finite decidable test.
For names in $T$ and $T'$, pass to
$T\otimes_S T'$ with residue map
$x\otimes y\mapsto\rho(x)\rho'(y)$ and use (C.3.HD7)
on their difference.  This decides equality even
when the two presented algebras are not domains.
Sums and products are computed in the same tensor
presentation.  Residues are given by $\rho$.
A named element is a unit precisely when its residue
is nonzero; its inverse is named by adjoining its
inverse to $T$, another pointed \'etale localization.
Finite residue lifts can be found once by enumerating
$S$ and testing its supplied surjection to $\kappa$.

If a canonical name from the enumerated standard
list is desired for any internally computed element,
search that list and apply the equality test.
Coverage proves that the search terminates.
Fix natural-number codes for the raw valid finite
names.  They form a decidable subset of $\mathbf N$.
Retain a name if no smaller valid name has the same
value, using only finitely many equality tests.
This gives canonical finite names as a decidable
subset of $\mathbf N$.  It works whether $R$ is
finite or infinite; subsequent searches run through
raw codes and never request a nonexistent next
new element of a finite ring.

\emph{Actual kernels, by flatness.}
Let a finite matrix $A$ over $R$ be given by such
names.  Put all its coefficients in one tensor
neighborhood $T$ and denote the descended matrix
by $A_T$.  Compute a finite generating set
$u_1,\ldots,u_c$ of $\ker A_T$ over $T$ by the
polynomial kernel algorithm.
Flatness of $T\to R$ in (C.3.HD6) gives
\[
            \ker A=(\ker A_T)\otimes_T R.
\]
Thus the images of the $u_j$ are generators of
the full actual kernel over $R$.  No approximation
or formal solution space is substituted here.
The generators need not be a free basis; none
of the classification procedures require one.

\emph{Application to representations and forms.}
The ring $R$ is Noetherian and henselian with finite
residue $\kappa$
\cite[Lemma 15.46.3, Tag 06LJ]{C-StacksHenselianEffective}.
The preceding paragraph supplies every kernel
calculation used in the proof of
Theorem~\ref{C:effective-henselian-DVR-WG}; this is
exactly where that proof used the DVR Smith procedure.
All other arguments there apply without change:
finite projectives over the local ring are free;
actual Hom generators have a computable finite
residue image; an intertwiner is invertible exactly
when its determinant has nonzero residue; and
Lemma~\ref{C:henselian-finite-algebra-lifting} supplies
semiperfectness and exact radical isometry lifting.
Endomorphism enumeration uses the computed finite
generators and all raw finite coefficient names,
so its positive lifting searches terminate even
when $R$ is finite or has zero divisors.

For a split idempotent, choose a basis of its
image by independent residue columns.  A unit
minor proves their independence, and Nakayama
proves generation.  This replaces any use of
Smith reduction to choose a free basis.
The least-code construction of the indecomposable
index set uses finitely many predecessor tests
and remains a decidable subset of $\mathbf N$.
Every residue algebra and its form data are
finite and computable as before.  The same actual
generators give (C.3.HD2), and the same tensor,
minor and finite coordinate algorithms give
all the claimed operations.  This proves the
extension with all its stated input access supplied.
\end{proof}
